% Mechanism derivation made before inspecting the new scientific outcomes.
\section{A Finite-Budget Ceiling on State-Selection Value}
An optimal conditional-risk rule does not by itself imply a practically large improvement.
In the scalar Gaussian continuation model one can bound the entire value of using the current noise state, before collecting scientific trajectories.

Fix the same paid probe, remaining budgets, public regime prior and finite catalogue for both rules.
Let $J_{j,a}^{r}$ be the unconditional commitment risk at initialization radius $r$ and let $s_j^2=\operatorname{Var}_j(m_j(Y))$.
Assume the common-state transfer coefficient $k_a$ is the same under the two variance hypotheses, as it is for the frozen learners with common persistence.
The even conditional risk score is
\[
L_a(y)=F_a(y)+k_a^2S(y),\quad
F_a(y)=\sum_jp_j(y)J_{j,a}^{r},\quad
S(y)=\sum_jp_j(y)[m_j(y)^2-s_j^2].
\]
Let $c(y)$ minimize $F_a(y)$ and let $u(y)$ minimize $L_a(y)$, using fixed measurable tie rules.
Both are even selectors and therefore have the initialization-minimax interpretation already proved.
In this section $R(v)$ means the maximum prior-weighted risk over the public scalar initialization interval; evenness attains that maximum at $e=\pm r$.
Then
\begin{equation}
0\le R(c)-R(u)
\le4\varphi(1)\bigl(\max_a k_a^2-\min_a k_a^2\bigr)
\sum_j\pi_js_j^2.
\label{eq:state-value-ceiling}
\end{equation}
The result is a finite-model upper bound on state information beyond the model posterior; it does not bound the value of identifying the covariance regime itself.
It remains valid for an observable mean-invariant Gaussian probe, using that probe's own explained variances.

For the proof, pointwise optimality gives
\[
0\le L_c-L_u\le (k_c^2-k_u^2)S
\le(\max_a k_a^2-\min_a k_a^2)|S|.
\]
The first inequality uses $L_u\le L_c$, and the second uses $F_c\le F_u$.
Since $p_j(y)g(y)=\pi_jf_j(y)$,
\[
\mathbb E_g|S(Y)|\le\sum_j\pi_j\mathbb E_j|m_j(Y)^2-s_j^2|
=4\varphi(1)\sum_j\pi_js_j^2.
\]
The last equality follows from the centered Gaussian identity $\mathbb E|G^2-1|=4\varphi(1)$; it also holds at $s_j=0$ by continuity.
Integrating proves Eq.~\eqref{eq:state-value-ceiling}.
The stop action has $k=0$ and fits the same bound.

For a running mean with $N$ applied observations and sample persistence $\rho=\lambda^b<1$,
\[
|k_a|=\frac{\rho(1-\rho^N)}{(N+1)(1-\rho)}
\le\frac{\rho}{(N+1)(1-\rho)}.
\]
For an averaged stable affine-delay learner, assume its impulse response is absolutely summable with sum of absolute values at most $G_*$.
The exact averaging weights obey $|w_t|\le\eta G_*/N$, giving
\[
|k_a|\le\frac{\eta G_*\rho}{N(1-\rho)}.
\]
Absolute summability is required here; it is not inferred merely from a numerical finite-horizon check.
For the registered affine steps it follows from the constant-delay characteristic polynomial $z^{D+1}-z^D+\alpha$, with $\alpha=\eta\mu=0.1/(D+1)$.
All roots are inside the unit disk for $0<\alpha<2\sin[\pi/(4D+2)]$.
Indeed for positive small $\alpha$ the root initially at one moves inside while the other roots remain near zero.
At a unit-circle crossing $z=e^{i\omega}$ the equation $z^D(z-1)=-\alpha$ gives a smallest positive crossing magnitude $2\sin[\pi/(4D+2)]$; root continuity excludes an earlier exit.
The registered $\alpha$ is strictly below that magnitude, since $2\sin[\pi/(4D+2)]\ge2/(2D+1)>0.1/(D+1)$.
The impulse generating function $[1-x+\alpha x^{D+1}]^{-1}$ is therefore analytic on a disk of radius greater than one, proving absolute summability.
The finite registered delay set has a finite uniform $G_*$.
If non-stop catalogue horizons are at least a positive constant times the budget scale $B$, and persistence and impulse bounds are uniform over that catalogue, the state-selection value is $O(B^{-2})$.
Relative to a best fixed risk of order $B^{-1}$ this contribution is at most $O(B^{-1})$.
For these stable averaged learners, with finite participation, nonzero private variance, persistence bounded below one and affordable horizons proportional to $B$, that fixed-risk order follows directly: $\sum_tw_t$ tends to a nonzero constant, Cauchy--Schwarz gives a private-variance lower bound proportional to $N^{-1}$, and the weight bound and summable temporal covariance give the matching upper bound.
Covariance-regime adaptation can change the leading $B^{-1}$ variance coefficient, so these two information roles must not be conflated.
This ceiling explains why a mathematically correct selection-dependence refinement can have little policy-level effect at long horizons, without converting that explanation into a claimed empirical success.
